% Figure: the KKT active-set example. Minimise (x-2)^2 + (y-2)^2 subject to
% x + y <= 1 and x >= 0. The unconstrained minimum (2,2) is infeasible; the
% constrained optimum is the foot of the perpendicular from (2,2) to the line
% x + y = 1, at (1/2, 1/2). Circular contours of f are centred on (2,2).
\begin{figure}[htbp]
\centering
\begin{tikzpicture}[scale=1.6,
  contour/.style={engblue, thick},
  cons/.style={engamber, very thick},
  feas/.style={englime!25},
  perp/.style={engred, thick, dashed},
]
  % Feasible region: x >= 0 and x + y <= 1, clipped to a viewing box.
  \begin{scope}
    \clip (-0.4, -0.9) rectangle (2.6, 2.6);
    \fill[feas] (0, -0.9) -- (0, 1) -- (1.9, -0.9) -- cycle;
  \end{scope}
  % Contours of f = (x-2)^2 + (y-2)^2, circles centred on (2,2).
  \foreach \r in {0.5, 1.0, 1.5, 2.121} {
    \draw[contour] (2, 2) circle (\r);
  }
  % Constraint lines.
  \draw[cons] (-0.4, 1.4) -- (1.9, -0.9) node[anchor=north west, pos=0.15]{$x + y = 1$};
  \draw[cons] (0, -0.9) -- (0, 2.6) node[anchor=west, pos=0.95]{$x = 0$};
  % Axes.
  \draw[engblack!45, -{Latex[length=1.4mm]}] (-0.4, 0) -- (2.6, 0) node[right]{$x$};
  \draw[engblack!45, -{Latex[length=1.4mm]}] (0, -0.9) -- (0, 2.6) node[above]{$y$};
  % Unconstrained minimum (infeasible).
  \filldraw[enggray] (2, 2) circle (0.04);
  \node[enggray, anchor=south west] at (2, 2) {$(2,2)$ infeasible};
  % Perpendicular from (2,2) to the line, foot at (0.5, 0.5).
  \draw[perp] (2, 2) -- (0.5, 0.5);
  \filldraw[engred] (0.5, 0.5) circle (0.05);
  \node[engred, anchor=north east] at (0.5, 0.5) {$(\tfrac12,\tfrac12)$};
\end{tikzpicture}
\caption[KKT active-set geometry.]{The two-inequality example of the
active-set sweep: minimise $(x - 2)^{2} + (y - 2)^{2}$ over the feasible
region $x \geq 0$, $x + y \leq 1$ (green). The objective's contours are
circles centred on the unconstrained minimum $(2, 2)$, which lies outside
the feasible region. The constrained optimum sits at the foot of the
perpendicular from $(2, 2)$ to the active line $x + y = 1$, at
$(\tfrac{1}{2}, \tfrac{1}{2})$ (red), where the smallest feasible contour
is tangent to the constraint. Only the first inequality is active there;
the case sweep discards the other three active-set assumptions on a
negative multiplier or a primal-feasibility violation.}
\label{fig:vol02:ch11:kkt-activeset}
\end{figure}
